\clearpage\thispagestyle{empty}
\pdfbookmark[0]{Sources for Gauss’s September 1802 Letter to Bode}{bode18020910-sources}
\begin{center}\textbf{Sources for Gauss’s September 1802 Letter to Bode}\end{center}
\noindent This AI translation follows Carl Friedrich Gauss’s letter to Johann Elert Bode of 10 September 1802 as printed in \href{https://gdz.sub.uni-goettingen.de/id/PPN236007467}{Gauss’s \textit{Werke}, vol.~VI} (Göttingen, 1874), p.~223.
\par\smallskip\noindent The \href{https://api.digitale-sammlungen.de/iiif/image/v2/bsb10538334_00237/full/1450,/0/default.jpg}{original printing} is on pp.~227--229 of Bode’s \href{https://wwwadd.zah.uni-heidelberg.de/publikationen/baj/baj1805/baj1805.htm}{\textit{Astronomisches Jahrbuch für das Jahr 1805}} (Berlin, 1802). The letter gives the elements used for the Pallas ephemeris, discusses the planet’s brightness and includes a note comparing the predictions of elements III. The \href{https://api.digitale-sammlungen.de/iiif/image/v2/bsb10538334_00238/full/1450,/0/default.jpg}{footnote on p.~228} is signed G. in the yearbook.
\par\smallskip\noindent The yearbook prints the ephemeris, with the \href{https://api.digitale-sammlungen.de/iiif/image/v2/bsb10538334_00239/full/1450,/0/default.jpg}{opposition times on p.~229}. The collected edition gives a bracketed description of the ephemeris and retains the two opposition lines. Both translations follow this collected scope and retain the unitless fields of the inclination row.
\par\smallskip\noindent Translation and typesetting produced by AI. Yearbook scans from the Bayerische Staatsbibliothek.
\clearpage
